\documentclass[10pt]{article}

\usepackage[utf8]{inputenc}

\usepackage[T1]{fontenc}

\usepackage{amsmath}

\usepackage{amsfonts}

\usepackage{amssymb}

\usepackage[version=4]{mhchem}

\usepackage{stmaryrd}



\begin{document}

ласть, заполненная жидкостью. На границе $\partial \Omega$ этой области обычно задаются условия прилипания:



\[

u(t, x)=0 \quad \text { при } \quad x=\left(x^{1}, x^{2}, x^{3}\right) \in \partial \Omega,

\]



а при $t=0$ начальное условие:



\[

u(0, x)=u_{0}(x) .

\]



Пусть $u_{0}(x)=u_{0}(x, \omega)$ - случайное векторное поле с распределением вероятностей $\mu$, а $u(t, x, \omega)$ - решение указанной задачи. Тогда мера $P$, задаюшая распределение вероятностей для $u(t, x, \omega)$, называется статистическим реш е н и е м, отвечающим начальной мере $\mu$.



М. И. Вишик, А. В. Фурсиков.



ГИДРОМАГНИТНАЯ ТУРБУЛЕНТНОСТЬ, магнит ногидродинамическая турбулентность,- совокупность случайных взаимодействующих магнитных и гидродинамических полей электропроводящей жидкости. Г. т. возбуждается и поддерживается в космич. условиях (ядра планет, межпланетная среда, конвективные зоны звезд, газовые диски в двойных звездных системах и галактиках), а также в нек-рых современных крупных технич. установках, напр. в жидкометаллич. теплоносителе ядерных реакторов-размножителей.



Обшей теории Г.т. пока не сушествует, ввиду математич. трудностей, стоящих при рассмотрении полной системы уравнений магнитной гидродинамики. Известен целый ряд частных, приближенных и модельных решений, а также результаты физич. экспериментов и космич. наблюдений. В одном предельном случае, когда можно пренебречь влиянием гидродинамич. движений на (сильное) магнитное поле, Г.т. преврашается в гидродинамич. турбулентность проводящей жидкости в заданном магнитном поле. В другом предельном случае, когда пренебрегают воздействием магнитного поля на движения, дело сводится к изучению поведения магнитного поля при заданном турбулентном течении (задачи гидромагнитного динамо). Оба предельных случая хорошо изучены и рассматриваются отдельно.



Полезной моделью служит двумерная Г. т. несжимаемой жидкости. Идеально проводяшее невязкое плоское течение характеризуется тремя квадратичными сохраняющимися величинами



\[

\int\left(\frac{\rho v^{2}}{2}+\frac{H^{2}}{8 \pi}\right) d x d y, \quad \int A^{2} d x d y, \quad \int \boldsymbol{v} \boldsymbol{H} d x d y

\]



где $\boldsymbol{v}, \boldsymbol{\rho}-$ скорость и плотность жидкости, $\boldsymbol{H}$ и $\boldsymbol{A}-$ магнитное поле и его вектор-потенциал. Первый интеграл - это полная (кинетическая плюс магнитная) энергия, третий называется кросс - п и ральностью, второй не имеет названия и возникает только в двумерном случае. В отличие от случая двумерной гидродинамич. турбулентности, в двумерной Г. т. энстрофия, то есть $\int(\operatorname{rot} v)^{2} d x d y$ не сохраняется. Если к двумерной диссипативной Г.т. подводить только кинетич. энергию, то, при условии убывания магнитного поля на бесконечности, магнитная энергия временно растет, а затем исчезает (Я. Б. Зельдович, 1956). Современные вычислительные эксперименты показали, что при инжектировании полной энергии возникает каскад энергии от больших к меньшим масштабам и обратный каскад величины $\int A^{2} d x d y$. Кроме того, при малой диссипации и достаточно плавном распределении скорости пространственные распределения вихря, плотности тока и магнитного поля оказываются существенно неоднородными, перемежаемыми, в то время как в двумерной гидродинамич. турбулентности распределение вихря достаточно однородно.



В трехмерном случае известна модель (Р. Крейкнан [2]), в к-рой Г.т. представляется совокупностью случайных несфазированных магнитногидродинамич. волн, распространяющихся в различных направлениях с альфвеновской ско- ростью в каждом масштабе. Отсюда вытекает равенство кинетич. и магнитной спектральной плотностей в каждом масштабе так наз. инерционного интервала. В слабой Г.т., когда преобладают тройные взаимодействия, это приводит к спектру $k^{-3 / 2}$. В этой модели, однако, весьма сушественно предположение о несфазированности волн, так как при этом теряется свойство перемежаемости Г.т., то есть присутствие структурных элементов, отмеченное выше на примере двумерной модели. Исследование Г. т. в приближении гидромагнитного динамо указывает на перемежаемость распределения магнитного поля, к-рое концентрируется в (сфазированные) жгуты или слои.



Особую роль в Г. т. могут играть топологич. особенности течения и магнитного поля, в частности присутствие средней спиральности (см. Спиральная турбулентность).



С появлением суперкомпьютеров начались численные исследования трехмерной Г.т.; напр., выполнено прямое интегрирование уравнений индукции и Навье - Стокса для несжимаемого потока со случайной короткокоррелированной силой, достигнув пространственного разрешения, соответствующего числам Рейнольдса порядка 100 (группа Ю. Фриша, см. [3]).



Лит.: [1] М офф а т Г. К., Возбуждение магнитного поля в проводяшей среде, пер. с англ., М., 1980; [2] Kraich n an R. H., «Physics Fluids», 1965, v. 8, p. 1385-87; [3] Meneguzzi M., FrischU., Pouquet A., «Phys. Rev. Lett.», 1981, v. 47, p. 1060-64.



А. А. Рузмайкин.



ГИДРОМАГНИТНОЕ ДИНАМО - см. Магнитной гидродинамики уравнения.



ГИЛЬБЕРТА - ШМИДТА OПЕРАТОР - см. Интегральный оператор.



ГИЛЬБЕРТА - ШМИДТА ТЕОРИЯ - см. Интегральное уравнение.



ГИЛЬБЕРТОВА АЛГЕБРА - алгебра $А$ с инволюцией над полем комплексных чисел, снабженная невырожденным скалярным произведением (|), причем выполняются следующие аксиомы:



\begin{enumerate}

  \item $(x \mid y)=\left(y^{*} \mid x^{*}\right)$ для всех $x, y \in A$;



  \item $(x y \mid z)=\left(y \mid x^{*} z\right)$ для всех $x, y, z \in A$;



  \item для всех $x \in A$ отображение $y \rightarrow x y$ пространства $A$ в $A$ непрерывно;



  \item множество элементов вида $x y, x, y \in A$, всюду плотно в $A$.



\end{enumerate}



П р имерам и Г.а. являются алгебры $L^{2}(G)$ (относительно свертки), где $G-$ компактная топологич. группа, и алгебра операторов Гильберта - Шмидта в данном гильбертовом пространстве.



Пусть $A$ - Г. а., $H$ - гильбертово пространство - пополнение $A, U_{x}$ и $V_{x}$ - элементы алгебры ограниченных линейных операторов в $H$, являющиеся продолжениями по непрерывности умножений слева и справа на $x$ в $A$. Отображение $x \rightarrow U_{x}$ (соответственно $x \rightarrow V_{x}$ ) есть невырожденное представление алгебры $A$ (соответственно алгебры с инволюцией, противоположной $A$ ) в гильбертовом пространстве $H$. Слабое замыкание семейства операторов $U_{x}$ (соответственно $V$ ) является а л ге бр о й $\mathrm{H}$ е й ма на в $H$; она называется л е в о й (соответственно п р а в о й) алгеброй Неймана данной Г.а. $A$ и обозначается $U(A)$ (соответственно $V(A)) ; U(A)$ и $V(A)$ являются коммутантами друг друга; это - полуконечные алгебры Неймана. Любая Г.а. однозначно определяет нек-рый точный нормальный полуконечный след на алгебре Неймана $U(A)$; обратно, если даны алгебра Неймана $\mathfrak{u}$ и точный нормальный полуконечный след на $\mathfrak{u}$, то можно построить Г. а. такую, что левая алгебра Неймана этой Г. а. изоморфна $\mathfrak{u}$, и след, определяемый на $\mathfrak{u}$ Г. а., совпадает с исходным (см. [1]). Таким образом, Г.а. является средством изучения полуконечных алгебр Неймана и следов на них; нек-рое обобщение понятия Г. а. позволяет изучать аналогичными средствами не обязательно полуконечные алгебры Неймана (см. [2|).



Jum.: 11] D i x mier J., Les algèbres d'opérateurs dans l'espace hilbertien, 2 ed., P., 1969; [2] Takesaki M., Tomita's theory of modular Hilbert algebras and its applications, B., 1970.



А. И. Штерн.





\end{document}